% !PRINTED 644
\weberpage{644}
%<<CONT
may in part be zero. By the same alteration, \S.\,185 (6) then yields
 the value of $\eta_0$.

With this, then, the restrictive assumptions of \S.\,185 are
removed, and we have the general theorem:

\gesperrt{I. Every root $\xi$ of a metacyclic equation
of prime degree $n$ can be represented in the
form}
\weq{9}{\xi = A + \sum_{0,\,n-2}^{\nu} K_\nu \tau_\nu^{r_{n-2}} \tau_{\nu+1}^{r_{n-3}} \cdots \tau_{\nu+n-2}^{r_0},}
\gesperrt{where $A$ is a rational quantity, $k_0, k_1 \dots k_{n-2}$ are the
roots, different from one another and from zero, of a cyclic
 equation of degree $(n-1)$, and $K_\nu$ is a rational function
of $k_\nu$, whose form is the same for all $\nu$. The exponents
$r_0, r_1 \dots r_{n-2}$ are the least positive residues of the
numbers $1, g, g^2 \dots g^{n-2}$, where $g$ is a primitive
root of $n$. The $n$ values that one obtains from (9)
when one assigns to the radicals $\tau_\nu = \sqrt[n]{k_\nu}$ their various
 values are the $n$ roots
of one and the same rational equation.}

The formula (9) results from \S.\,185 (6) when $A$ and $K_\nu$
are replaced by $nA$ and $nK_\nu$, which has been done for simplification.

But the converse of Theorem I can now also be stated.

To prove this, we denote by $\varepsilon$ any
imaginary $n$th root of unity and set
\weq{10}{\begin{aligned}\xi_h &= A + \sum_{0,\,n-2}^{\nu} \varepsilon^{h r_\nu} K_\nu \tau_\nu^{r_{n-2}} \tau_{\nu+1}^{r_{n-3}} \cdots \tau_{\nu+n-2}^{r_0}\\ h &= 0, 1, 2 \dots n-1,\end{aligned}}
or, abbreviated,
\weq{11}{\xi_h = A + \sum_{0,\,n-2}^{\nu} \varepsilon^{h r_\nu} K_\nu \sqrt[n]{R_\nu},}
where, as before,
\weq{12}{\sqrt[n]{R_\nu} = \tau_\nu^{r_{n-2}} \tau_{\nu+1}^{r_{n-3}} \cdots \tau_{\nu+n-2}^{r_0}}
is set. We now have to investigate the changes
that result for $\xi_h$ when we

1. determine one of the radicals $\tau_\nu = \sqrt[n]{k_\nu}$ differently, and

2. permute the $k_0, k_1 \dots k_{n-2}$ cyclically.
